\section*{Capítulo \ref{ch:b3:bioinformatics} --- Bioinformática e análise de sequências}

\begin{solution}{exo:b3:bioinformatics:1}
Global: as duas sequências alinhadas de ponta a ponta, cada resíduo
numa coluna --- para duas proteínas que se acreditam homólogas em todo
o comprimento, como as ortólogas de uma enzima de manutenção. Local: o
par de subcadeias de maior escore, ignorado o resto --- para achar um
domínio compartilhado (um domínio SH2 em duas proteínas de sinalização
de resto não aparentadas) ou um gene numa longa sequência genômica.
\end{solution}

\begin{solution}{exo:b3:bioinformatics:2}
Bordas $0,-1,-2,-3$ nos dois sentidos. Linha A: $1, 0, -1$. Linha G:
$0, 0, -1$. Linha C: $-1, -1, 1$. Ótimo $F(3,3) = 1$: AGC sobre AAC sem
lacunas (correspondência, discordância, correspondência:
$1 - 1 + 1 = 1$).
\end{solution}

\begin{solution}{exo:b3:bioinformatics:3}
$s(a,a) = \lambda^{-1}\log\bigl(q_{aa}/p_{a}^{2}\bigr)$. O triptofano é
raro ($p_{W} \approx 0.013$), de modo que a chance de dois triptofanos
se alinharem ao acaso, $p_{W}^{2}$, é ínfima, e um par de triptofanos
conservado é um sinal de homologia muito mais forte que um par de
leucinas conservado ($p_{L}
\approx 0.1$); a razão de log-chances é
correspondentemente maior.
\end{solution}

\begin{solution}{exo:b3:bioinformatics:4}
O \omterm{thm:b3:bioinformatics:evalue}{valor E} é o número de \omterm{def:b3:bioinformatics:alignment}{alinhamentos} com escore ao menos tão alto que
se esperaria por acaso numa busca desta consulta contra um banco deste
tamanho. $E = 3$ significa que três escores desses são esperados por
acaso: o acerto não é evidência de homologia (ele ainda pode sê-lo, mas
a busca não tem como dizer).
\end{solution}

\begin{solution}{exo:b3:bioinformatics:5}
$mn = 400\times 2\times 10^{11} = 8\times 10^{13} = 2^{46.2}$. $E(45) =
2^{1.2} \approx 2.3$; $E(55) = 2^{-8.8} \approx 2\times 10^{-3}$;
$E(65) = 2^{-18.8} \approx 2\times 10^{-6}$. $E = 10^{-3}$ exige
$S' =
46.2 + 10.0 = 56$ bits. Uma consulta de \num{40} resíduos tem um
$mn$ dez vezes menor, $2^{42.9}$: $53$ bits bastam --- mas uma consulta
curta raramente chega sequer a isso.
\end{solution}

\begin{solution}{exo:b3:bioinformatics:6}
\omterm{def:b3:bioinformatics:motif}{Conteúdos de informação}: $2$, $1$, $0$ e $2 - H$ com $H = -(0.7\log_{2}
0.7 + 3\times 0.1\log_{2} 0.1) = 0.36 + 1.00 = 1.36$, logo $0.64$.
Total $R = 3.64$ bits. Correspondências ao acaso: $9.2\times 10^{6}$
posições em duas fitas $\times 2^{-3.64} \approx 7\times 10^{5}$ --- o
\omterm{def:b3:bioinformatics:motif}{motivo} é quase inútil sozinho.
\end{solution}

\begin{solution}{exo:b3:bioinformatics:7}
Uma inserção de vários resíduos é um único evento mutacional, de modo
que seu custo não deveria crescer linearmente com seu comprimento: um
custo de abertura mais um pequeno custo de extensão modela isso. Uma
penalidade de abertura alta demais força discordâncias onde caberia uma
lacuna e desalinha tudo depois de uma inserção verdadeira; uma
penalidade baixa demais espalha lacunas por toda parte, emparelhando
resíduos ao acaso e inflando a identidade.
\end{solution}

\begin{solution}{exo:b3:bioinformatics:8}
Não necessariamente: o \omterm{def:b3:bioinformatics:alignment}{alinhamento} cobre um segmento de \num{130}
resíduos, que é a assinatura de um domínio compartilhado, e não de uma
ortóloga alinhada em todo o seu comprimento. Teste: busque a proteína
da mosca de volta contra o proteoma humano (a consulta é seu melhor
acerto, em todo o comprimento?), identifique o domínio com um HMM de
\omterm{def:b3:bioinformatics:msa}{perfil} e construa uma árvore de genes da família em várias espécies.
\end{solution}

\begin{solution}{exo:b3:bioinformatics:9}
Um \omterm{def:b3:bioinformatics:msa}{perfil} registra, coluna por coluna, o que a família tolera: uma
posição que é sempre hidrofóbica mas nunca o mesmo resíduo, um resíduo
catalítico invariante, uma posição que é sempre uma lacuna em metade da
família. Um \omterm{def:b3:bioinformatics:alignment}{alinhamento} par a par pontua cada resíduo contra um único
outro resíduo e não tem como saber que uma valina na posição 40 é ``tão
boa quanto'' a isoleucina que ali está na consulta. O \omterm{def:b3:bioinformatics:msa}{perfil} também
pondera as colunas conservadas, de modo que uma similaridade fraca
concentrada onde a família é conservada se torna significativa.
\end{solution}

\begin{solution}{exo:b3:bioinformatics:10}
Com \omterm{prop:b3:bioinformatics:logodds}{escore esperado} positivo $\mu > 0$ por coluna, o escore acumulado
ao longo da diagonal de duas sequências aleatórias é um passeio
aleatório com deriva positiva: após $n$ colunas ele vale cerca de
$\mu n$, de modo que o melhor \omterm{def:b3:bioinformatics:alignment}{alinhamento local} é essencialmente a
coisa toda e seu escore cresce como $\mu n$, e não como $\log n$. A
teoria de Karlin--Altschul, que exige deriva negativa para que escores
altos sejam excursões raras, não se aplica e nenhum $\lambda$ existe.
Para DNA a \qty{60}{\%} de GC, a chance de correspondência é
$2(0.3^{2}) + 2(0.2^{2}) = 0.26$, de modo que o \omterm{prop:b3:bioinformatics:logodds}{escore esperado} é
$0.26 - 0.74 = -0.48$: ainda negativo, e a estatística vale; mas um
esquema como correspondência $+1$, discordância $-0.3$ teria esperança
$+0.04$ e reportaria o genoma inteiro como um só \omterm{def:b3:bioinformatics:alignment}{alinhamento}.
\end{solution}

\begin{solution}{exo:b3:bioinformatics:11}
\omterm{def:b3:bioinformatics:blast}{Sementes} e encadeamento: encontre correspondências exatas ou quase
exatas de $k$-meros entre as duas sequências com uma tabela de
dispersão, guarde as que se alinham em diagonais coerentes, encadeie-as
e rode a \omterm{thm:b3:bioinformatics:nw}{programação dinâmica} só nas lacunas entre as \omterm{def:b3:bioinformatics:blast}{sementes}
encadeadas; abre mão de \omterm{def:b3:bioinformatics:alignment}{alinhamentos} em regiões sem \omterm{def:b3:bioinformatics:blast}{semente} (trechos
muito divergentes). Faixas: se as duas sequências são sabidamente quase
colineares, calcule apenas as células dentro de uma faixa de largura
$w$ em torno da diagonal, a um custo $wn$ em vez de $mn$; abre mão de
qualquer \omterm{def:b3:bioinformatics:alignment}{alinhamento} com uma inserção maior que a faixa.
\end{solution}

\begin{solution}{exo:b3:bioinformatics:12}
A sequência codificadora tem período três: as três posições do códon
têm composições de bases diferentes (a terceira é a mais enviesada), e
o uso de códons é desigual em cada espécie. Um modelo com três estados
codificadores em sequência, cada um emitindo bases com a composição
daquela posição do códon, atribui ao DNA codificante uma probabilidade
maior do que o estado não codificante atribui, ao longo de uma janela
de algumas dezenas de códons, mesmo sem os códons de parada. Num genoma
humano os éxons são curtos (\qty{150}{bp}), separados por íntrons de
quilobases, de modo que o sinal codificante é breve e interrompido; o
modelo também precisa reconhecer os sítios de splicing, que são sinais
fracos, e a enorme quantidade de sequência não codificante produz
muitos segmentos codificantes falsos.
\end{solution}

\begin{solution}{pb:b3:bioinformatics:1}
\textbf{1.} Linhas K, Q, T; colunas K, A, Q, T; bordas $0,-1,-2,-3,-4$
e $0,-1,-2,-3$. Linha K: $1, 0, -1, -2$; linha Q: $0, 0, 1, 0$; linha
T: $-1, -1, 0, 2$. Ótimo $2$: \texttt{K-QT} sobre \texttt{KAQT}.
\textbf{2.} Melhor escore local $3$: \texttt{CAT} contra \texttt{CAT}
(resíduos 4--6 de GATCAT com 2--4 de ACAT); \texttt{ATCAT} contra
\texttt{A-CAT} também pontua $4 - 1 = 3$.
\textbf{3.} $300\times 450 = 1.35\times 10^{5}$ atualizações; contra o
banco, $300\times 1.2\times 10^{11} = 3.6\times 10^{13}$.
\textbf{4.} $3.6\times 10^{4}$ s, dez horas por consulta; as \omterm{def:b3:bioinformatics:blast}{sementes}
do \omterm{def:b3:bioinformatics:blast}{BLAST} saltam quase toda a tabela e respondem em segundos.
\textbf{5.} $q_{ab} = p_{a}p_{b}e^{\lambda s}$. Alanina: $e^{1.39} = 4.0$,
$q_{AA} = 0.074^{2}\times 4.0 = 0.022$, razão $4$. Triptofano:
$e^{3.82} = 45$, $q_{WW} = 0.013^{2}\times 45 = 0.0077$, razão $45$. Um
par de triptofanos alinhado é $45$ vezes mais frequente em homólogas do
que por acaso, e um par de alaninas apenas quatro vezes; ainda assim os
pares de alanina são mais comuns em termos absolutos, porque a alanina
é comum.
\textbf{6.} \qty{24}{\%} está na zona crepuscular, em que os
\omterm{def:b3:bioinformatics:alignment}{alinhamentos} aleatórios chegam a \qtyrange{15}{20}{\%}; o \omterm{thm:b3:bioinformatics:evalue}{valor E} do
\omterm{def:b3:bioinformatics:alignment}{alinhamento}, \omterm{def:b3:bioinformatics:motif}{motivos} conservados nas posições certas, uma
correspondência a um HMM de \omterm{def:b3:bioinformatics:msa}{perfil} de família conhecida ou um
enovelamento compartilhado resolveriam a questão.
\textbf{7.} $mn = 300\times 1.2\times 10^{11} = 3.6\times 10^{13}$;
$\log_{2}(mn) = 45.0$.
\textbf{8.} $E(92) = 2^{45 - 92} = 2^{-47} \approx 7\times 10^{-15}$;
$E(38) = 2^{7} = 128$.
\textbf{9.} $E = 10^{-3}$ em $S' = 45 + 10 = 55$ bits; $E = 1$ em $45$
bits.
\textbf{10.} Dez vezes o $mn$: $E \approx 7\times 10^{-14}$, ainda
esmagador.
\textbf{11.} Com $E = 128$, o décimo acerto é o que o acaso produz; uma
identidade de \qty{40}{\%} ao longo de $60$ resíduos não é evidência.
Uma busca por \omterm{def:b3:bioinformatics:msa}{perfil} dos resíduos 200--260 contra o banco de domínios
poderia mostrar se aquele segmento é um domínio conhecido, com uma
estatística que falta à comparação par a par.
\textbf{12.} Melhores acertos recíprocos em todo o comprimento são
compatíveis com uma ortologia um a um; não a provam --- uma duplicação
numa das linhagens depois da separação dá dois \omterm{def:b3:genomics:comparative}{co-ortólogos}, e a perda
da ortóloga verdadeira pode deixar uma paráloga como melhor acerto. Uma
árvore de genes com várias espécies é o teste.
\textbf{13.} $R = 2 + 2 + 1.6 + 2 + 0.8 + 1.2 + 0.4 + 0.3 = 10.3$ bits.
\textbf{14.} $3.2\times 10^{9}\times 2^{-10.3} \approx 2.5\times 10^{6}$
correspondências ao acaso.
\textbf{15.} $\log_{2}(3.2\times 10^{9}/200) = \log_{2}(1.6\times 10^{7})
\approx 24$ bits.
\textbf{16.} A coocorrência a espaçamento fixo soma os bits: $10.3 + 8 =
18.3$, o que supre $8$ dos $13.7$ que faltam; cerca de $5.7$ bits (um
fator de $50$ nas correspondências ao acaso) precisam vir de outro
lugar --- acessibilidade da \omterm{def:b3:chromatin-epigenetics:nucleosome}{cromatina}, outros parceiros.
\textbf{17.} $H = -(0.5\log_{2}0.5 + 0.5\log_{2}0.5) = 1$ bit; $R = 2 -
1 = 1$ bit.
\textbf{18.} Um fator bacteriano precisa achar seus poucos sítios num
genoma de \qty{4.6}{Mb} sem ajuda alguma da \omterm{def:b3:chromatin-epigenetics:nucleosome}{cromatina}: ele precisa de
uns $19$ bits, e seus sítios são longos e conservados. Um genoma
eucarionte é mil vezes maior e exigiria dez bits a mais, e ainda assim
seus fatores têm sítios curtos; eles alcançam especificidade por
combinação e pela restrição da \omterm{def:b3:chromatin-epigenetics:nucleosome}{cromatina} acessível, o que também torna
a regulação mais evoluível, já que um sítio curto é facilmente ganho ou
perdido.
\textbf{19.} $3/64 = 0.047$; o número esperado de códons antes de um de
parada é $64/3 \approx 21$.
\textbf{20.} $p_{A} = p_{T} = 0.31$, $p_{G} = p_{C} = 0.19$: $P(\text{TAA})
= 0.31^{3} = 0.030$, $P(\text{TAG}) = P(\text{TGA}) = 0.31^{2}\times 0.19
= 0.018$; total $0.066$, comprimento esperado de fase $15$ códons. O DNA
rico em AT é cheio de códons de parada, de modo que as fases abertas
aleatórias são mais curtas e as longas se destacam mais.
\textbf{21.} $(61/64)^{100} = e^{-4.80} = 0.008$; $(61/64)^{300} =
e^{-14.4} = 5.6\times 10^{-7}$.
\textbf{22.} Cada fase aberta maximal termina num códon de parada, e
$3.2\times
10^{9}$ inícios de códon contêm $3.2\times 10^{9}\times 3/64 = 1.5\times
10^{8}$ códons de parada: cerca de
$1.5\times 10^{8}\times 0.008 = 1.2\times 10^{6}$ fases aleatórias de
ao menos $100$ códons, e $1.5\times 10^{8}\times
5.6\times 10^{-7} \approx 80$ de ao menos $300$.
\textbf{23.} Uma bactéria de \qty{4.6}{Mb} tem uns $4\times 10^{5}$
códons de parada e, portanto, cerca de \num{3500} fases de acaso de
$100$ códons, mas quase nenhuma de $300$; seus genes têm em média $300$
códons e \qty{88}{\%} do DNA é codificante, de modo que uma fase aberta
longa é quase sempre um gene. No verme, \qty{1.5}{\%} do DNA codifica,
os éxons têm em média $50$ códons --- menos que o limiar do acaso --- e
um milhão de fases aleatórias de $100$ códons os soterra. Os
localizadores de genes eucariontes usam sinais de sítio de splicing,
viés de códons num modelo oculto de Markov, homologia com proteínas
conhecidas e, acima de tudo, transcritos sequenciados.
\textbf{24.} Uma leitura de um mensageiro processado se alinha ao genoma
em dois pedaços separados por um íntron: a divisão marca os dois sítios
de splicing base a base; a cobertura de leituras delineia os éxons e as
leituras pareadas ligam éxons sucessivos num mesmo transcrito,
resolvendo qual entre vários sítios candidatos de splicing é usado.
\textbf{25.} $E \approx 7\times 10^{-15}$ para o melhor acerto; cerca de
$24$ bits para especificar $200$ sítios no genoma; umas $80$ fases de
leitura de acaso de $300$ códons no genoma inteiro.
\end{solution}
